\chapter{Loss of Appetite}

\section*{Definition}
Reduced desire to eat prolonged anorexia can lead unintended weight loss nutritional deficiencies.

\section*{What Happens in Your Body}
Distinction from nausea (urge to vomit) and early satiety (feeling full quickly); central hypothalamic regulation peripheral gastrointestinal signals interact influencing appetite.

\section*{Causes}
\begin{itemize}
  \item Depression/anxiety
  \item Chronic infections
  \item Cancer/malignancy
  \item Thyroid disorders
  \item Medication side effects
\end{itemize}

\section*{When to See a Doctor}
See your doctor if:
\begin{itemize}
  \item You have unintentional weight loss along with loss of appetite
\end{itemize}

\section*{Related Symptoms}
\begin{itemize}
  \item Weight loss
  \item Fatigue
  \item Nausea
  \item Taste changes
  \item Fullness after small meals
\end{itemize}


\section*{Self-Care / Home Management}
\begin{itemize}
  \item Maintain fluids and try small, frequent meals containing protein and energy-dense foods that are tolerated.
  \item Address reversible contributors such as dental problems, nausea, pain, constipation, medication effects, depression, or difficulty obtaining food with professional help.
  \item Keep a record of intake and weight; seek medical assessment for persistent loss of appetite, dehydration, swallowing difficulty, or unintentional weight loss.
\end{itemize}

\section*{How Your Doctor Will Evaluate You}
\begin{itemize}
  \item History assesses weight change, intake, mood, swallowing, gastrointestinal symptoms, systemic symptoms, medicines, substance use, and food access.
  \item Examination includes nutrition, hydration, oral health, abdominal findings, and signs of endocrine, infectious, malignant, or psychiatric disease.
  \item Blood tests and targeted imaging or endoscopy are guided by the history, examination, age, and alarm features.
\end{itemize}

\section*{Treatment Options}
\begin{itemize}
  \item Treatment targets the cause and may include medication review, dental or gastrointestinal treatment, nutritional support, psychotherapy, or treatment of an underlying systemic illness.
  \item Appetite-stimulating medicines are not routine and should be considered only by a clinician after risks and benefits are reviewed.
\end{itemize}

\section*{References}

\begin{thebibliography}{99}
\bibitem{loss_of_appetite:harrison-appetite} Longo DL, et al. \emph{Harrison's Principles of Internal Medicine}. 21st ed. McGraw-Hill; 2022. Chapter 46: Anorexia and Cachexia.
\bibitem{loss_of_appetite:uptodate-anorexia} Evans WJ, et al. Anorexia and cachexia in adults. In: UpToDate, Post TW (Ed), UpToDate, Waltham, MA; 2025.
\bibitem{loss_of_appetite:schwartz} Schwartz MW, et al. Central regulation of appetite. \emph{Nature}. 2023;620(7973):274-285.
\bibitem{loss_of_appetite:feuerstein} Feuerstein JD, et al. Evaluation of unintentional weight loss. \emph{JAMA}. 2024;331(10):864-874.
\bibitem{loss_of_appetite:walston} Walston JD, et al. Anorexia of aging: pathophysiology and management. \emph{Lancet Healthy Longev}. 2023;4(6):e281-e291.
\bibitem{loss_of_appetite:morley} Morley JE, et al. \emph{The Endocrinology of Aging}. 5th ed. Springer; 2024. Chapter 8: Anorexia of Aging.
\end{thebibliography}
